\documentclass[../main.tex]{subfiles}
\begin{document}
%==========================================TIÊU ĐỀ TẬP========================================
\begin{table}[h]
  \begin{adjustwidth}{-1cm}{}
    \begin{tabular}{>{\centering\arraybackslash}p{6cm}|>{\raggedright\arraybackslash}p{12.5cm}}
      \multirow{5}{6cm}
      % Cột bên trái: ÉPISODE và số tập
      {\\[-40pt]
        \flaregothic\fontsize{20pt}{20pt}\selectfont \centering
        \color{eptype}ÉPISODE\color{black}\\
        \burbank\fontsize{72pt}{72pt}\selectfont
        \color{epnum}9\color{black}\\[6pt]
        \cafeteria\fontsize{20pt}{20pt}\selectfont\color{epnum}(1-5)\color{black}
        \\[18pt]
        \flaregothic\fontsize{18pt}{18pt}\selectfont
        \color{epattr}ÉPISODE PERDU
        \color{black}
      }
       &
      % Dòng thứ nhất: tên tập tiếng Nhật
      {\fotskip\fontsize{18pt}{18pt}\selectfont そのあと\ruby{墜}{つい}\ruby{落}{らく}した【\ruby{全}{ぜん}\ruby{米}{べい}が\ruby{泣}{な}いた】
      } \\[6pt]
      % Dòng thứ hai: tên tập tiếng Anh
       & {\cafeteria\fontsize{18pt}{18pt}\selectfont Dying for Hiko-boshi}                                            \\[4pt]
      % Dòng thứ ba: tên tập tiếng Việt
       & {\vnmsans\fontsize{18pt}{18pt}\selectfont Aqua và ngày lễ Thất Tịch}                                         \\[8pt]
      % Dòng thứ tư: Nhân vật xuất hiện
       & {\caviar{\fontsize{14pt}{14pt}\selectfont Nhân vật: \emp{kuon1} \emp{ryuuhou1} \emp{fubuki} \emp{aqua} \emp{mio}}} \\[8pt]
      % Dòng cuối cùng: Thời gian diễn ra của tập (thường sẽ là ngày phát hành)
       & {\continuum\fontsize{16pt}{16pt}\selectfont Ngày phát hành: 07/07/2019}                                      \\[18pt]
    \end{tabular}
  \end{adjustwidth}
\end{table}
%=========================================TRUYỆN CHÍNH========================================
\color{black}

\txtbx[2]{
    {\color{vol} \uline{Tóm tắt:}} Vào ngày lễ Thất Tịch, Minato Aqua khao khát được bay lượn trên bầu trời như trong phim hoạt hình. Nhờ sự trợ giúp bằng phương pháp đầy điên rồ của cậu Tổng Kuon, cô Hầu gái đã bay vút lên không trung nhưng lại không biết cách dừng lại và trôi dạt tận vào vũ trụ bao la. Tại đây, để giúp đỡ Ngưu Lang (Fubuki) vượt qua dải Ngân Hà bị lũ lụt để gặp gỡ Chức Nữ (Mio), Aqua đã hy sinh hai bím tóc của mình tạo ra lỗ đen liên kết và rơi tự do xuống Trái Đất. Tập phim khép lại với chương ngoại truyện ấm cúng nhưng không kém phần hài hước tại gia đình Hikawa, nơi Kuon bị mẹ và em gái trêu chọc khi thấy Aqua xuất hiện trên bản tin thời sự từ tận rừng Amazon.
}

Ngày mùng bảy tháng bảy, bầu trời Nhật Bản trong vắt, điểm xuyết vài dải mây trắng mỏng manh lững lờ trôi. Hôm nay là ngày lễ Thất Tịch, dịp để mọi người gửi gắm tâm tình và những điều ước giản dị lên các mảnh giấy tanzaku sắc màu, treo chúng lên cành tre xanh với hy vọng chúng sẽ thành hiện thực. Trong khi bầu không khí lễ hội ngập tràn trên khắp các ngả đường, văn phòng của \hll\ vẫn rộn ràng tiếng cười nói và những hoạt động thường nhật mang màu sắc độc nhất vũ trụ.

\img{1-5a}

Tại một góc nhỏ trong văn phòng, Aqua đang ngồi bó gối trên sàn nhà, mắt dán chặt vào chương trình hoạt hình chiếu trên chiếc màn hình lớn. Trên đó, một cậu bé đang tự do bay lượn giữa tầng mây nhờ một chiếc chong chóng kỳ lạ gắn trên đầu\footnote{Một liên hệ đến bảo bối \textit{Chong chóng tre} trong Doraemon.}. Đôi mắt cô hầu gái sáng bừng lên lấp lánh, tràn ngập sự ngưỡng mộ và khao khát trẻ con: ``Mình cũng muốn được bay lượn tự do như thế quá đi\dots''

Kuon ngồi trên chiếc ghế sofa đối diện, chứng kiến cảnh đó chỉ biết thở dài ngao ngán. Cậu thầm nghĩ trong lòng: ``\textit{Làm sao mà bay được cơ chứ\dots\ Nhưng mà nếu nói thẳng ra thì tội nghiệp em ấy, mất cả hứng khởi.}''

Thế nhưng, sau khi ngẫm nghĩ lại về những sự kiện kỳ quặc liên tiếp xảy ra tại văn phòng này suốt thời gian qua, cậu lại thay đổi ý nghĩ: ``\textit{Cơ mà cái văn phòng hololive này vốn dĩ đã bất bình thường rồi, chuyện quái gì mà chẳng xảy ra được. Hay là mình thử giúp em ấy một tay xem sao nhỉ?}''

Cậu chống cằm suy ngẫm một lúc, rồi một ý tưởng táo bạo nhưng không kém phần điên rồ lóe lên trong đầu. Cậu nghiêng người về phía Aqua, mỉm cười: ``Anh có thể giúp em thực hiện điều ước đó đấy.''

Aqua lập tức quay phắt lại, đôi mắt mở to sáng lấp lánh đầy kinh ngạc lẫn hy vọng ngập tràn: ``Thật ạ!? Anh làm cách nào thế ạ?''

Kuon khẽ gật đầu, nở một nụ cười nửa tự tin nửa e ngại: ``Ừ, thật chứ. Nhưng nói trước là phương pháp này nghe có hơi điên rồ một chút đấy\dots''

Cậu bắt đầu thì thầm giải thích cho Aqua cách thức để có thể cất cánh bay lên: một phương pháp nghe qua thì đầy tính hoang đường nhưng lại vô cùng hợp lý dưới góc nhìn khoa học giả tưởng của cậu Tổng. Còn cụ thể đó là gì, thì có lẽ chỉ có bộ não dị biệt của Kuon mới nghĩ ra được.

\begin{center}
  \uline{Ngày hôm sau.}
\end{center}

Trải qua vô số lần tập luyện vất vả, Aqua rốt cuộc đã thực hiện được ước mơ cháy bỏng của mình. Cô bay vút lên trời cao, lơ lửng giữa khoảng không lộng gió với hai bím tóc màu hành tím đang xoay tít như cánh quạt động cơ. Cô reo lên đầy phấn khích và sung sướng: ``Ha ha! Mình làm được thật rồi này!''

\gif{1-5b}{1}{24}

Nhưng niềm vui ngắn chẳng tày gang, chỉ vài giây sau, nụ cười trên môi nàng Hầu gái lập tức cứng đờ. Aqua bắt đầu mất kiểm soát, đôi mắt xoay như chong chóng đầy hoảng loạn: ``Ơ\dots\ nhưng mà làm thế nào để dừng lại và hạ cánh bây giờ!? Kuon-senpai vẫn chưa kịp chỉ cho mình cách quay về thế nào mà!?''

Aqua hoảng hốt bay loạn xạ trên bầu trời, tạo nên một màn trình diễn nhào lộn dở khóc dở cười nhưng cũng vô cùng nguy hiểm. Vừa chao lượn vô định, cô vừa gào thét thất thanh: ``\textbf{Ai đó làm ơn cứu tôi với! Cứu tôi xuống mau đi!}''

Cứ thế, hai bím tóc của cô cứ như hai cánh quạt trực thăng chạy hết công suất, liên tục đưa cô hướng thẳng lên trên cao. Cô vượt qua những tầng mây trắng bồng bềnh, đâm xuyên qua bầu khí quyển Trái Đất và cuối cùng trôi dạt vào tận vũ trụ bao la lạnh lẽo. Giữa không gian lấp lánh ánh sao huyền ảo, Aqua bất ngờ chạm trán một vị thần đang cưỡi trên lưng trâu. Điều khiến cô sững sờ nhất chính là ngoại hình của vị thần này trông quen thuộc đến kỳ lạ.

\img{1-5c}

 ``Ể!? Fubuki-chan đấy à!?'' Aqua trợn tròn mắt kinh ngạc.

Vị thần nọ bật cười khà khà, lắc đầu phủ nhận: ``Fubuki-chan nào chứ? Ta chính là sao Ngưu Lang (Hikoboshi) hiển hách đây mà!''

Aqua ngượng ngùng gãi đầu cười trừ: ``Ồ\dots\ hóa ra là con nhìn nhầm người. Con xin lỗi ạ\dots''

Vị thần cáo trắng khẽ thở dài đầy u sầu, bắt đầu giãi bày nỗi lòng tâm sự của mình: ``Hôm nay tuy là ngày lễ Thất Tịch, nhưng do trời đổ mưa lớn đột ngột nên dải Ngân Hà đã bị ngập lụt nghiêm trọng. Ta không có cách nào băng qua dòng nước lũ để gặp gỡ người yêu của mình là sao Chức Nữ (Orihime) được nữa rồi\footnote{Câu chuyện của thần Hikoboshi-Fubuki là sự tích ngày lễ Thất tịch ở Nhật Bản. Bắt nguồn từ truyền thuyết về chuyện tình của Chức Nữ (nàng tiên) và Ngưu Lang (chăn bò). Họ rất yêu nhau, nhưng bị Thiên đình phản đối do sự khác biệt về địa vị; Chức Nữ thuộc tiên giới, còn Ngưu Lang là người phàm. Họ rời xa nhau va mỗi người ở một bên dải ngân hà. Cách duy nhất để họ có thể gặp nhau là đúng vào ngày Thất tịch (7/7 Âm lịch), một con chim ác là sẽ tạo một cây cầu băng qua hai ngân hà để đưa họ với nhau, nhưng họ chỉ có thể gặp nhau mỗi năm một lần.}\dots''

Aqua đứng nghe mà mồ hôi hột tuôn ra như tắm, chỉ biết nở một nụ cười gượng gạo. Cô thầm nghĩ trong lòng: ``\textit{Mình không ngờ thần tiên mà cũng có nhiều tâm sự thế này đấy\dots\ Cơ mà sao cái cốt truyện này nghe quen thuộc giống hệt như trong mấy cuốn sách truyền thuyết thế nhỉ?}''

Cô định cất tiếng thắc mắc, nhưng khi nhìn vào đôi mắt đượm buồn và gương mặt đầy vẻ u sầu của vị thần chăn bò, Aqua lại cảm thấy lòng mình nhói lên. Cô quyết định gạt phắt mọi nghi vấn sang một bên, nắm chặt hai tay hạ quyết tâm: ``Thần đừng quá lo lắng! Con sẽ tạo ra một lỗ đen vũ trụ để thần có thể băng qua dải Ngân Hà và hội ngộ với cô ấy ngay lập tức!''

Vị thần cáo trắng nghe thấy thế liền giật bắn mình, hai tai cáo dựng đứng lên như bị điện giật. Ông cuống quýt vung tay phản đối kịch liệt: ``Con điên rồi hả!? Nếu làm chuyện đó, con sẽ phải hy sinh hai bím tóc động cơ của mình cho Ngục giới, có nghĩa là chúng sẽ vĩnh viễn biến mất khỏi thế gian này đấy!''

Ống kính lập tức quay cận cảnh vào hai bím tóc của Aqua, lúc này vẫn đang quay phần phật với tốc độ chóng mặt như hai cánh quạt phản lực.
 
``Chưa kể, nếu mất đi hai bím tóc đó, con sẽ hoàn toàn mất khả năng bay lượn và rơi tự do thẳng xuống Trái Đất đấy!'' vị thần ra sức ngăn cản và thuyết phục cô từ bỏ ý định liều lĩnh. Ông thậm chí còn đưa ra một hình ảnh minh họa đầy bi kịch: cảnh cô nằm sấp dưới một hố đất sâu, tư thế tử vong kinh điển giống hệt một nhân vật có tên bắt đầu bằng chữ Y trong bộ anime có số bảy\footnote{Ý nói đến cảnh cái chết của Yamcha trong \textit{Dragon Ball Z} (\textit{Bảy viên ngọc rồng}).}.

\begin{figure}[H]
  \centering
  \includegraphics[width=12.5cm]{../../../../Image/Book1/1-5d.png}
  \caption{Hình ảnh Aqua sẽ hy sinh, giống như ``ai đó bắt đầu bằng chữ Y trong bộ anime có số 7 trong đó'' theo tưởng tượng đầy bi kịch của vị thần.}
\end{figure}

Thế nhưng, nàng Hầu gái bướng bỉnh vẫn kiên định với quyết định của mình. Đôi mắt cô rực sáng lên niềm tin mãnh liệt: ``Chuyện nguy hiểm hay mất tóc gì đó con không quan tâm!''

Vị thần Hikoboshi hoàn toàn chấn kinh, đôi tai cụp xuống rồi lại dựng lên, ánh mắt trợn tròn trước sự quả cảm đến mức điên rồ của cô hầu gái nhỏ. Aqua dõng dạc hét lớn, giọng điệu đầy vẻ hào hùng: ``Con là hầu gái của \hll\ mà! Sứ mệnh của con chính là mang lại hạnh phúc cho tất cả mọi người, dù có phải hy sinh bản thân đi chăng nữa!''

Lời tuyên bố đanh thép vang vọng khắp khoảng không vũ trụ tĩnh mịch, tựa như một khúc tráng ca tôn vinh hành động cao cả của cô. Không chần chừ thêm một giây, Aqua giơ tay giật mạnh hai bím tóc màu hành tím của mình, xoay tít chúng như hai bánh quay ly tâm rồi dùng hết sức bình sinh ném thẳng ra khoảng không vô định.

\img{1-5e}

Ngay lập tức, một lỗ đen khổng lồ được hình thành, xoáy tròn cuồn cuộn với nguồn năng lượng dữ dội hút cạn ánh sáng xung quanh. Sức hút kinh hoàng kéo tuột cả thần Hikoboshi cùng chú trâu của ông vào bên trong. Giữa cảnh tượng hỗn loạn đó, tiếng la thất thanh của vị thần vọng lại: ``\textbf{Con đang làm cái trò điên rồ gì thế này!?}''

Aqua không đáp lại. Thể lực cạn kiệt cùng việc mất đi hai bím tóc động cơ khiến cô không còn lực nâng để duy trì vị trí. Cô bắt đầu mất kiểm soát và rơi tự do xuống dưới. Thời gian như ngừng trôi, cảnh vật chuyển động chậm rãi: Aqua rơi rụng giữa bầu trời đêm đầy sao lấp lánh, trông cô lúc này kiêu hãnh như một vị anh hùng vừa hoàn thành sứ mệnh cứu rỗi thế giới.

\img{1-5f}

Khẽ nhắm nghiền đôi mắt, cô hầu gái mỉm cười thanh thản, sẵn sàng đón nhận kết cục tồi tệ nhất. Ở phía xa, tiếng hét tuyệt vọng của vị thần Ngưu Lang vẫn văng vẳng: ``Aqua! Đừng làm thế! Nếu con rơi xuống thì--''

Nhưng mọi chuyện đã quá muộn. Aqua rơi vút qua các tầng khí quyển của Trái Đất, để lại một vệt đuôi rực rỡ như một ngôi sao băng quả cảm. Trong khoảng không vũ trụ, lỗ đen ma thuật vẫn tiếp tục xoay vần, và nhờ có nó, thần Hikoboshi rốt cuộc đã vượt qua được dải Ngân Hà ngập lụt để xuất hiện bên cạnh người yêu của mình.

Thế nhưng, ở một nơi khác dưới mặt đất\dots

\ct

Cơn mưa rào mùa hạ đang trút xuống thành phố mỗi lúc một nặng hạt.

Bên trong văn phòng hololive tĩnh lặng, Ookami Mio đứng tựa bên cửa sổ, ánh mắt đượm buồn dõi theo những giọt nước mưa lăn dài trên mặt kính loang lổ. Tiếng mưa rơi tí tách đơn điệu ngoài kia dường như càng làm tăng thêm vẻ trầm lắng của căn phòng. Cô khẽ thở dài, đưa ngón tay vẽ những vòng tròn vô định lên lớp sương mờ bám trên mặt kính và lẩm nhẩm: ``Chán thật đó\dots\ Hôm nay đã là lễ Thất Tịch rồi mà trời lại mưa tầm tã thế này\dots''

Đứng ngay cạnh cô, Shirakami Fubuki cũng đang chăm chú nhìn ra màn mưa xám xịt bên ngoài. Cô gật đầu đồng cảm, khẽ đáp lại bằng một giọng điệu có chút xa xăm: ``Đúng thật nhỉ\dots''

Mio khẽ cụp tai xuống, giọng cô thoảng nhẹ như gió thoảng: ``Bà có nghĩ năm nay Chức Nữ và Ngưu Lang vẫn có thể gặp nhau không? Mưa lớn thế này dải Ngân Hà chắc ngập lụt mất rồi, chắc là họ không gặp được nhau đâu nhỉ\dots''

Fubuki đột nhiên quay sang nhìn người bạn sói của mình, khóe môi khẽ cong lên thành một nụ cười ấm áp. Cô cắt ngang dòng suy nghĩ u sầu của Mio bằng một câu khẳng định đầy chắc chắn: ``Họ sẽ gặp được nhau thôi. Tôi biết chắc chắn là như thế mà\footnote{Ở đây, Chức Nữ đại diện cho Mio (Orihime) và Ngưu Lang đại diện cho Fubuki (Hikoboshi), phản ánh mối quan hệ vô cùng gắn bó giữa hai người.}.''

Mio ngơ ngác quay sang, đôi mắt ngập tràn vẻ ngạc nhiên nhìn thẳng vào Fubuki. Trong khoảnh khắc ngắn ngủi ấy, cô dường như thoáng thấy điều gì đó vừa quen thuộc vừa kỳ lạ trong ánh mắt của cô bạn cáo trắng, tựa như một sợi dây liên kết vô hình kéo dài qua ngàn năm ký ức bị lãng quên. ``Fubuki\dots\ Sao bà lại biết chắc chắn như vậy?''

\img{1-5g}

Fubuki chỉ khẽ mỉm cười đầy ẩn ý, đôi tai cáo mềm mại khẽ rung lên nhè nhẹ theo nhịp điệu của tiếng mưa rơi. Cô không giải thích thêm, chỉ quay đầu tiếp tục ngắm nhìn thế giới mờ ảo qua làn nước. Những giọt mưa rơi xuống không ngừng, nối liền đất trời bằng những sợi tơ bạc mỏng manh.

Giai điệu nhẹ nhàng từ bài hát \textit{Rain} của Motohiro Hata phát ra từ chiếc loa nhỏ ở góc văn phòng, len lỏi vào từng kẽ hở của không gian, kể lại một câu chuyện tình yêu vượt qua dòng chảy thời gian. Sự lựa chọn nhạc nền này thật sự vô cùng hòa hợp với bầu không khí lãng mạn pha chút trầm buồn của ngày lễ hôm nay.

Mio không hỏi thêm điều gì nữa. Cô lặng lẽ đứng sát lại gần Fubuki, cùng chia sẻ một khoảnh khắc yên bình, sâu lắng giữa tiết trời giông bão. Cơn mưa bên ngoài vẫn không có dấu hiệu dứt hạt\dots\ thế nhưng trong tâm tư của hai cô gái, dường như một cây cầu Ô Thước vô hình đã được bắc qua, kết nối những nhịp đập cảm xúc không cần diễn tả bằng lời.

%==========================================TRUYỆN PHỤ=========================================
\omk

\begin{center}
  {\textit{*Các nhân vật nói bằng tiếng Etsunan.}}
\end{center}

\pov{Hikawa Kuon}

Hôm nay là ngày lễ Thất Tịch, và cả gia đình tôi quyết định ở nhà dưỡng sức. Một phần vì cơn mưa mùa hạ đang trút xuống như trút nước ngoài kia, phần còn lại thì - nói thật lòng - tôi đang ế chỏng ế chơ. Chẳng ai trong nhà có chút hứng thú nào với việc ra ngoài trong thời tiết ẩm ướt khó chịu này.

Không phải vì tôi bất tài đến mức không kiếm nổi mảnh tình vắt vai, chỉ là tôi thực sự không có thời gian và tâm trí cho chuyện đó. Công việc quản lý tại \hll\ phải được đặt lên hàng đầu, còn chuyện yêu đương cứ để tùy duyên đi.

Mẹ tôi vừa bưng những bát chè đậu đỏ mát lạnh từ dưới bếp lên chia cho cả nhà. Đây là phong tục ăn chè đậu đỏ cầu duyên của ngày Thất Tịch, và với tôi, nó giống như một lời nhắc nhở ngọt ngào nhưng đầy sát thương rằng tôi vẫn đang cô đơn lẻ bóng. Nhưng thôi, ít nhất thì tôi vẫn được thưởng thức đồ ngọt miễn phí.

\chrint

Nói thêm một chút về mẹ tôi, bà Hikawa Kaien (\ruby{火}{ひ}\ruby{川}{かわ}\ruby{海}{かい}\ruby{英}{えん}). Tên thời con gái của bà là Musada Kaien (\ruby{無}{む}\ruby{佐}{さ}\ruby{田}{だ}\ - \ruby{Mư}{Vô}\ruby{xa}{Tá}\ruby{đa}{Điền}\ruby{海}{かい}\ruby{英}{えん}). Do luật pháp ở Etsunan khác biệt với Nhật Bản, cho phép người phụ nữ sau khi kết hôn được giữ lại họ thời con gái của mình, nên trong sổ hộ khẩu gia đình tôi, tên bà được ghi song song dưới dạng Hikawa Kaien / Musada Kaien . Mẹ tôi là một người phụ nữ cực kỳ tự do, phóng khoáng và có phần ngang bướng, chẳng chịu để ai sai bảo dễ dàng. 

Thời còn trẻ, bà từng là một MC kiêm ca sĩ hát hội chợ có tiếng. Khoảng thời gian đi diễn đó, bà dành phần lớn thời gian ở bên ngoài nhiều hơn ở nhà. Đó là lý do vì sao tôi và em gái từ nhỏ đã gắn bó và dành nhiều thời gian bên bố hơn là bên mẹ. Bản thân bố tôi cũng rất tôn trọng sự tự do của bà nên luôn tạo điều kiện để bà theo đuổi đam mê. Phải mãi cho đến khi bà về hưu, đồng thời những bất hòa, hiểu lầm trong gia đình được giải quyết triệt để, bà mới chịu lui về chăm lo chu toàn cho tổ ấm. Hiện tại, ngoài công việc nội trợ, bà còn đảm nhận việc quản lý kho vật liệu xây dựng cho em trai bà (tức cậu ruột tôi). 

Về trang phục thường ngày, mẹ tôi rất giản dị. Bà chỉ thích mặc áo phông hoặc áo dài tay kết hợp với quần đùi, quần soóc hoặc quần dài tùy theo mùa. Những bộ đầm lộng lẫy chỉ xuất hiện mỗi khi bà đi diễn ca hát hội chợ ngày xưa mà thôi.

\chrint

Cô em gái đang cặm cụi ăn chè đối diện tôi tên là Hikawa Ryuuhou (\ruby{火}{ひ}\ruby{川}{かわ}\ruby{劉}{りゅう}\ruby{宝}{ほう}), tên đầy đủ trên giấy tờ của nó là Hikawa Musada Ryuuhou (mang cả họ bố lẫn họ mẹ). Nó kém tôi mười tuổi, sở hữu đôi mắt hai màu cực kỳ đặc biệt thừa hưởng từ bố mẹ: mắt phải màu đỏ của bố và mắt trái màu xanh lam thẫm của mẹ. 

Ryuuhou có mái tóc cắt bob dài ngang cổ và thường để xõa, hiếm khi lắm mới cột lên để cho gọn. Trái ngược với vẻ ngoài có phần tomboy tinh nghịch, thường ngày nó chỉ thích mặc áo phông hoặc áo dài tay kết hợp với quần soóc hoặc quần dài năng động. Con bé hầu như không bao giờ chịu mặc váy, trừ khi có những sự kiện đặc biệt, như là đi sự kiện hay đi biểu diễn. Nó luôn viện cớ là sợ ông anh trai này nhìn trộm - dẫu tôi thề có trời đất chứng giám là tôi chẳng bao giờ thèm làm cái trò biến thái đó cả! Thực chất, lý do chính chỉ là do tính cách tomboy phóng khoáng, ghét gò bó của nó mà thôi. 

Con bé này hiện đang là trưởng nhóm của một nhóm nhạc underground và đã tự kiếm được một khoản tiền khá khẩm, thậm chí còn nhiều hơn cả thu nhập của bố mẹ tôi thời điểm đó. Tôi từng có một khoảng thời gian làm quản lý bất đắc dĩ cho nhóm nhạc của nó. Một phần vì bị đứa em gái mè nheo nhờ vả, phần khác vì tôi muốn xem cái nhóm nhạc ``bát nháo'' toàn những cô cậu nhóc kia có làm nên trò trống gì không. Nhưng nhờ tính cách hướng ngoại bẩm sinh cùng giọng hát cực tốt được cả bố lẫn mẹ rèn giũa từ bé, nhóm nhạc của Ryuuhou dần đi vào quỹ đạo ổn định chỉ sau khoảng ba tháng và bắt đầu gặt hái những quả ngọt đầu tiên. Sau khi thấy mọi chuyện đã ổn thỏa, tôi chủ động rút lui để nhóm tự quản lý. Cũng chính nhờ kinh nghiệm làm quản lý bất đắc dĩ cho nhóm nhạc của em gái mà tôi đã tích lũy được những kiến thức nền tảng về công việc của một Tổng quản, giúp tôi vượt qua buổi phỏng vấn xin việc đầy cam go với A-chan và YAGOO ở \hll\ sau này.

Tuy nhiên, cái gì cũng có giá của nó. Vì quá mải mê tập luyện và đi diễn cùng nhóm nhạc trong suốt hai năm trời, sức học của Ryuuhou sa sút nghiêm trọng. Hậu quả là nó bị lưu ban mất hai năm, dẫn đến việc giờ này vẫn đang chuẩn bị lẹt đẹt bước vào lớp sáu thay vì lớp tám như bạn bè đồng trang lứa.

\chrint

Để xua bớt bầu không khí có phần ảm đạm của buổi chiều mưa, tôi bảo bố bật tivi lên xem tin tức. Thế nhưng, ngay khi bản tin thời sự quốc tế vừa hiện ra, tôi suýt chút nữa đã phun hết ngụm chè đậu đỏ trong miệng ra sàn nhà.

Trên màn hình tivi, dòng chữ tiêu đề chạy ngang nổi bật: `Bắt giữ du khách nhập cảnh trái phép tại rừng rậm Amazon'. Cảnh quay chuyển đến một phóng viên đang che ô đứng dưới cơn mưa rừng tầm tã, giọng đọc dồn dập: ``Chính quyền sở tại vừa thông báo bắt giữ một nữ du khách người Nhật Bản có hành vi nhập cảnh bất hợp pháp vào sâu trong khu bảo tồn Amazon. Theo lời khai ban đầu của đối tượng\dots''

Khung hình ngay sau đó phóng to vào gương mặt của kẻ phạm tội. Dù nhà đài đã tinh tế làm mờ vùng mắt, nhưng làm sao tôi có thể nhìn nhầm được cơ chứ? Cái mái tóc tím củ hành cùng bộ trang phục hầu gái quen thuộc kia chắc chắn 100\% là Minato Aqua! Dưới ống kính máy quay của phóng viên quốc tế, cô hầu gái nhà tôi giơ ngón tay cái lên, nở một nụ cười rạng rỡ và tuyên bố đầy tự hào trước truyền thông thế giới: ``Ít nhất đôi uyên ương đó đã được hạnh phúc bên nhau, thế là quá ổn với con rồi!''

Meme ``Miễn là cô ấy hạnh phúc thì tôi cũng ổn' kinh điển được tái hiện hoàn hảo ngay giữa rừng rậm Nam Mỹ. Tôi đơ toàn tập, mắt chữ O miệng chữ A nhìn chằm chằm vào màn hình ti vi, đầu óc quay cuồng không hiểu nổi làm thế nào mà em ấy lại có thể bay từ Nhật Bản sang tận bên kia bán cầu chỉ bằng hai bím tóc.''

Đúng lúc đó, bố tôi liếc mắt sang tôi, vẻ mặt đầy tò mò: ``Bố nhớ cô bé tóc hồng-xanh này là đồng nghiệp cùng công ty của con đúng không, Kuon?''

Chương trình truyền hình ngay lập tức bị mẹ tôi chen ngang với nụ cười chọc ghẹo đầy ẩn ý: ``Ồ, bạn gia đình con đấy à? Hay là bạn gái của con thế?''

Tôi hoảng loạn xua tay loạn xạ, mặt đỏ bừng lên vì xấu hổ, cố gắng thanh minh trong vô vọng: ``Không phải đâu mẹ ơi! Oan cho con quá, con hoàn toàn không liên quan gì đến vụ vượt biên này hết!''

Nhưng như để đổ thêm dầu vào lửa, cô em gái Ryuuhou liền cất giọng hát bài đồng dao châm chọc đầy đắc ý: ``Onii-san biết yêu rồi\~{} Onii-san biết yêu rồi kìa\~{}!''

Mới nghe thế, mẹ tôi đã tủm tỉm cười, nhấp một ngụm trà: ``Mẹ thấy bát chè đậu đỏ hôm nay linh nghiệm thật đấy chứ.''

Cuối cùng, tôi chỉ biết ôm đầu gào lên trong sự bất lực tột cùng trước màn trêu chọc tập thể của gia đình: ``Đã bảo là không phải rồi mà!!!''

\end{document}
